% GENERATED FILE. Edit canonical lessons, then run npm run generate:books.
% canonical-source-hash: fnv1a64:da36bec8b9edc9a6
% canonical-lessons: GU-C498-viday, GU-C498-kanya, GU-C498-kunvaru, GU-C498-dahej, GU-C498-chutacheda

\chapter{The Bride Leaves Home}
\label{ch:gu-the-bride-leaves-home}
\hlchaptermodality{498}

\begin{chapteropening}
\textbf{By the end of this chapter:} \emph{I can talk about a bride's farewell.}

Complete the last lesson of chapter 498: I can talk about a bride's farewell.
\end{chapteropening}

\section[\texorpdfstring{\gu{વિદાય} (vidāy)}{વિદાય (vidāy)}]{\gu{વિદાય} (vidāy) — a farewell; the bride's send-off}

\label{lesson:GU-C498-viday}



\begin{quote}
Before the new one: say the Gujarati for the groom's procession, then the Gujarati for to marry off.
\end{quote}

\subsection*{You'll want to know: \gu{વિદાય}}
\textbf{\gu{વિદાય}} — \emph{vidāy} — \textquotedblleft{}a farewell; the bride's send-off\textquotedblright{}.

\begin{quote}
\gu{દીકરીની} \gu{વિદાય} \gu{વખતે} \gu{માતા} \gu{રડી} \gu{પડી}.
\end{quote}

\begin{quote}
\emph{dīkrīnī vidāy vakhte mātā raḍī paḍī.}
\end{quote}

\begin{quote}
At her daughter's send-off the mother burst into tears.
\end{quote}

A feminine noun.

\begin{cousinweb}
\textbf{\gu{વિદાય}} (\emph{vidāy}) is recorded by Monier-Williams as Sanskrit \emph{vidāya}, \textquotedblleft{}permission to go away, dismissal with good wishes\textquotedblright{}, with a note that in this sense it is perhaps not a Sanskrit word at all and a comparison with Arabic \emph{wadāʿ}, \textquotedblleft{}farewell\textquotedblright{}; its history is not settled.
\end{cousinweb}

\begin{grammarlens}[title={What you've built}]
One more word for describing what you lived and felt.
\end{grammarlens}

\subsection*{Your turn}
\begin{itemize}
\raggedright
  \item \emph{Say it:} \emph{vidāy}
  \item \emph{Say it:} \emph{vidāy}, once more
  \item \emph{Say it:} \emph{dīkrīnī vidāy vakhte mātā raḍī paḍī.} — At her daughter's send-off the mother burst into tears.
  \item \emph{Recall:} say the Gujarati for the groom's procession, then the Gujarati for to marry off, then say \emph{vidāy} again
\end{itemize}

\begin{tcolorbox}[breakable,colback=teal!4,colframe=teal!35!black,title={Before you move on}]
What does \gu{વિદાય} mean? (\textquotedblleft{}A farewell; the bride's send-off\textquotedblright{}.) Say it once more.
\end{tcolorbox}
\section[\texorpdfstring{\gu{કન્યા} (kanyā)}{કન્યા (kanyā)}]{\gu{કન્યા} (kanyā) — a bride; a girl}

\label{lesson:GU-C498-kanya}



\begin{quote}
Before the new one: say the Gujarati for to marry off, then the Gujarati for a farewell.
\end{quote}

\subsection*{You'll want to know: \gu{કન્યા}}
\textbf{\gu{કન્યા}} — \emph{kanyā} — \textquotedblleft{}a bride; a girl\textquotedblright{}.

\begin{quote}
\gu{લાલ} \gu{સાડી} \gu{પહેરીને} \gu{કન્યા} \gu{બહાર} \gu{આવી}.
\end{quote}

\begin{quote}
\emph{lāl sāḍī paherīne kanyā bahār āvī.}
\end{quote}

\begin{quote}
Dressed in a red sari, the bride came out.
\end{quote}

A feminine noun.

\begin{cousinweb}
\textbf{\gu{કન્યા}} (\emph{kanyā}) is a tatsama, Sanskrit \emph{kanyā}, \textquotedblleft{}a girl, a daughter\textquotedblright{}, as Monier-Williams records it.
\end{cousinweb}

\begin{grammarlens}[title={What you've built}]
One more word for describing what you lived and felt.
\end{grammarlens}

\subsection*{Your turn}
\begin{itemize}
\raggedright
  \item \emph{Say it:} \emph{kanyā}
  \item \emph{Say it:} \emph{kanyā}, once more
  \item \emph{Say it:} \emph{lāl sāḍī paherīne kanyā bahār āvī.} — Dressed in a red sari, the bride came out.
  \item \emph{Recall:} say the Gujarati for to marry off, then the Gujarati for a farewell, then say \emph{kanyā} again
\end{itemize}

\begin{tcolorbox}[breakable,colback=teal!4,colframe=teal!35!black,title={Before you move on}]
What does \gu{કન્યા} mean? (\textquotedblleft{}A bride; a girl\textquotedblright{}.) Say it once more.
\end{tcolorbox}
\section[\texorpdfstring{\gu{કુંવારું} (kunvārũ)}{કુંવારું (kunvārũ)}]{\gu{કુંવારું} (kunvārũ) — unmarried}

\label{lesson:GU-C498-kunvaru}



\begin{quote}
Before the new one: say the Gujarati for a farewell, then the Gujarati for a bride.
\end{quote}

\subsection*{You'll want to know: \gu{કુંવારું}}
\textbf{\gu{કુંવારું}} — \emph{kunvārũ} — \textquotedblleft{}unmarried\textquotedblright{}.

\begin{quote}
\gu{મારા} \gu{મામા} \gu{હજી} \gu{કુંવારા} \gu{છે}.
\end{quote}

\begin{quote}
\emph{mārā māmā hajī kunvārā che.}
\end{quote}

\begin{quote}
My uncle is still unmarried.
\end{quote}

\begin{cousinweb}
\textbf{\gu{કુંવારું}} (\emph{kunvārũ}) is inherited from Sanskrit \emph{kumāra}, \textquotedblleft{}a boy, a youth\textquotedblright{}; CDIAL lists the Gujarati word, \textquotedblleft{}unmarried\textquotedblright{} (CDIAL 3303).
\end{cousinweb}

\begin{grammarlens}[title={What you've built}]
One more word for describing what you lived and felt.
\end{grammarlens}

\subsection*{Your turn}
\begin{itemize}
\raggedright
  \item \emph{Say it:} \emph{kunvārũ}
  \item \emph{Say it:} \emph{kunvārũ}, once more
  \item \emph{Say it:} \emph{mārā māmā hajī kunvārā che.} — My uncle is still unmarried.
  \item \emph{Recall:} say the Gujarati for a farewell, then the Gujarati for a bride, then say \emph{kunvārũ} again
\end{itemize}

\begin{tcolorbox}[breakable,colback=teal!4,colframe=teal!35!black,title={Before you move on}]
What does \gu{કુંવારું} mean? (\textquotedblleft{}Unmarried\textquotedblright{}.) Say it once more.
\end{tcolorbox}
\section[\texorpdfstring{\gu{દહેજ} (dahej)}{દહેજ (dahej)}]{\gu{દહેજ} (dahej) — a dowry}

\label{lesson:GU-C498-dahej}



\begin{quote}
Before the new one: say the Gujarati for a bride, then the Gujarati for unmarried.
\end{quote}

\subsection*{You'll want to know: \gu{દહેજ}}
\textbf{\gu{દહેજ}} — \emph{dahej} — \textquotedblleft{}a dowry\textquotedblright{}.

\begin{quote}
\gu{દહેજ} \gu{માગવું} \gu{ખોટું} \gu{છે}.
\end{quote}

\begin{quote}
\emph{dahej māgvũ khoṭũ che.}
\end{quote}

\begin{quote}
Demanding a dowry is wrong.
\end{quote}

A neuter noun.

\begin{cousinweb}
\textbf{\gu{દહેજ}} (\emph{dahej}) is usually traced to Arabic \emph{jahāz}, \textquotedblleft{}equipment, a trousseau\textquotedblright{}, which reached Gujarati through Persian.
\end{cousinweb}

\begin{grammarlens}[title={What you've built}]
One more word for describing what you lived and felt.
\end{grammarlens}

\subsection*{Your turn}
\begin{itemize}
\raggedright
  \item \emph{Say it:} \emph{dahej}
  \item \emph{Say it:} \emph{dahej}, once more
  \item \emph{Say it:} \emph{dahej māgvũ khoṭũ che.} — Demanding a dowry is wrong.
  \item \emph{Recall:} say the Gujarati for a bride, then the Gujarati for unmarried, then say \emph{dahej} again
\end{itemize}

\begin{tcolorbox}[breakable,colback=teal!4,colframe=teal!35!black,title={Before you move on}]
What does \gu{દહેજ} mean? (\textquotedblleft{}A dowry\textquotedblright{}.) Say it once more.
\end{tcolorbox}
\section[\texorpdfstring{\gu{છૂટાછેડા} (chūṭācheḍā)}{છૂટાછેડા (chūṭācheḍā)}]{\gu{છૂટાછેડા} (chūṭācheḍā) — a divorce}

\label{lesson:GU-C498-chutacheda}



\begin{quote}
Before the new one: say the Gujarati for unmarried, then the Gujarati for a dowry.
\end{quote}

\subsection*{You'll want to know: \gu{છૂટાછેડા}}
\textbf{\gu{છૂટાછેડા}} — \emph{chūṭācheḍā} — \textquotedblleft{}a divorce\textquotedblright{}.

\begin{quote}
\gu{દસ} \gu{વરસ} \gu{પછી} \gu{તેમના} \gu{છૂટાછેડા} \gu{થયા}.
\end{quote}

\begin{quote}
\emph{das varas pachī temnā chūṭācheḍā thayā.}
\end{quote}

\begin{quote}
After ten years they got divorced.
\end{quote}

A masculine noun, used in the plural.

\begin{cousinweb}
\textbf{\gu{છૂટાછેડા}} (\emph{chūṭācheḍā}) is built from \emph{chūṭā}, \textquotedblleft{}loosened, apart\textquotedblright{}, on the stem of \textbf{\gu{છૂટવું}} (\emph{chūṭvũ}), \textquotedblleft{}to get loose\textquotedblright{} (CDIAL 3707), and \emph{cheḍā}, \textquotedblleft{}ends\textquotedblright{}: the ends of cloth tied together at a wedding come apart.
\end{cousinweb}

\begin{grammarlens}[title={What you've built}]
One more word for describing what you lived and felt.
\end{grammarlens}

\subsection*{Your turn}
\begin{itemize}
\raggedright
  \item \emph{Say it:} \emph{chūṭācheḍā}
  \item \emph{Say it:} \emph{chūṭācheḍā}, once more
  \item \emph{Say it:} \emph{das varas pachī temnā chūṭācheḍā thayā.} — After ten years they got divorced.
  \item \emph{Recall:} say the Gujarati for unmarried, then the Gujarati for a dowry, then say \emph{chūṭācheḍā} again
\end{itemize}

\begin{tcolorbox}[breakable,colback=teal!4,colframe=teal!35!black,title={Before you move on}]
What does \gu{છૂટાછેડા} mean? (\textquotedblleft{}A divorce\textquotedblright{}.) Say it once more.
\end{tcolorbox}
